\documentclass[11pt,a4paper]{article}

\usepackage{amsmath,amssymb,amsthm}
\usepackage{mathtools}
\usepackage{hyperref}
\usepackage{cleveref}
\usepackage[margin=1in]{geometry}
\usepackage{enumitem}
\usepackage{booktabs}

\newtheorem{theorem}{Theorem}[section]
\newtheorem{lemma}[theorem]{Lemma}
\newtheorem{proposition}[theorem]{Proposition}
\newtheorem{corollary}[theorem]{Corollary}
\newtheorem{definition}[theorem]{Definition}
\newtheorem{axiom}{Axiom}
\newtheorem{conjecture}[theorem]{Conjecture}
\newtheorem{remark}[theorem]{Remark}
\newtheorem{example}[theorem]{Example}

\newcommand{\Prob}{\mathsf{Prob}}
\newcommand{\ProbProp}{\mathsf{ProbProp}}
\newcommand{\E}{\mathbb{E}}
\newcommand{\pand}{\wedge_p}
\newcommand{\por}{\vee_p}
\newcommand{\pnot}{\neg_p}
\newcommand{\pzero}{\mathbb{0}}
\newcommand{\pone}{\mathbb{1}}
\newcommand{\ple}{\leq_p}
\newcommand{\plt}{<_p}
\newcommand{\padd}{+_p}
\newcommand{\pmul}{\cdot_p}
\newcommand{\psub}{-_p}
\newcommand{\Nat}{\mathbb{N}}
\newcommand{\Real}{\mathbb{R}}

\title{Synthetic Probability Type Theory\\[0.5em]
\large Towards Constructing Probability from Primitives}

\author{Probabilistic Foundations Project}

\date{\today}

\begin{document}

\maketitle

\begin{abstract}
We prove that probabilistic Zorn's lemma holds for finite types without any form of choice: every finite chain-complete probabilistic poset has an $\varepsilon$-near-maximal element. We then sketch how the infinite case follows from projective limits of finite approximations---constructing distributions rather than selecting points. Along the way, we reduce the axiomatic footprint of our probabilistic foundations from 48 to 43 axioms by observing that several axioms follow from commutativity, with further reductions possible given auxiliary lemmas. We conclude by posing open questions about synthetic probability type theory, where $\Prob$ would be a primitive type former.
\end{abstract}

\tableofcontents

% ============================================================================
\section{Introduction}
% ============================================================================

\subsection{Motivation}

Papers 1--4 developed probabilistic foundations using 48 axioms. This paper asks: how many of these axioms are truly independent? Can we construct $\Prob$ rather than axiomatize it?

\subsection{Goals}

\begin{enumerate}
    \item Reduce the axiom count by proving derivable axioms as theorems
    \item Prove probabilistic Zorn for finite types (no choice needed)
    \item Sketch the proof for infinite types via projective limits
    \item Outline a fully synthetic approach where probability is primitive
\end{enumerate}

\subsection{Summary of Results}

\begin{center}
\begin{tabular}{lccc}
\toprule
Result & Formalized & Documented & Axioms Affected \\
\midrule
Commutativity reductions & Yes (Lean) & \S\ref{sec:axiom-reduction} & $-5$ \\
De Morgan derivation & No & \S\ref{sec:axiom-reduction} & $-2$ (blocked) \\
CondExp reduction & Partial & \S\ref{sec:axiom-reduction} & $-3$ (needs work) \\
Finite prob\_zorn & No & \S\ref{sec:finite-zorn} & $+0$ (new theorem) \\
Infinite prob\_zorn & No & \S\ref{sec:projective-limits} & $-2$ (conjecture) \\
\bottomrule
\end{tabular}
\end{center}

% ============================================================================
\section{Part 1: Axiom Reduction}
\label{sec:axiom-reduction}
% ============================================================================

The probabilistic foundations of Papers 1--4 rest on 48 axioms. We systematically analyze which are independent and which are derivable.

\subsection{Current Axiom Inventory}

The axioms fall into six categories:

\begin{center}
\begin{tabular}{lrl}
\toprule
Category & Count & Examples \\
\midrule
$\Prob$ type and operations & 7 & $\Prob$, $\pzero$, $\pone$, $\ple$, $\padd$, $\pmul$, $\psub$ \\
Order structure & 5 & reflexivity, transitivity, antisymmetry, bounds \\
Multiplication & 7 & identity, commutativity, associativity, monotonicity \\
Addition & 4 & identity, commutativity, associativity \\
Subtraction & 7 & self-inverse, cancellation, complement \\
Distributivity \& misc & 4 & distribution, non-triviality \\
\midrule
$\Prob$ subtotal & 34 & \\
\midrule
Conditional expectation & 7 & normalization, monotonicity, product, etc. \\
Countable summation & 5 & $\mathsf{prob\_sum}$ and properties \\
Probabilistic Zorn & 2 & $\mathsf{prob\_zorn}$, $\mathsf{prob\_zorn\_witness}$ \\
\midrule
\textbf{Total} & \textbf{48} & \\
\bottomrule
\end{tabular}
\end{center}

\subsection{Derivable Axioms: Commutativity}

Five axioms are immediate consequences of commutativity:

\begin{proposition}[Trivial Reductions]
\label{prop:trivial-reductions}
The following five axioms are derivable from other axioms:\footnote{Each derivation is a one-line rewrite, formalized in \texttt{Synthetic.lean}.}
\begin{enumerate}
    \item $p \pmul \pone = p$ from $\pone \pmul p = p$ via $\pmul$-commutativity
    \item $p \padd \pzero = p$ from $\pzero \padd p = p$ via $\padd$-commutativity
    \item $p \pmul \pzero = \pzero$ from $\pzero \pmul p = \pzero$ via $\pmul$-commutativity
    \item $(p \padd q) \pmul r = (p \pmul r) \padd (q \pmul r)$ from $r \pmul (p \padd q) = (r \pmul p) \padd (r \pmul q)$ via $\pmul$-commutativity (right-distributivity from left-distributivity)
    \item $\pone \psub \pone = \pzero$ as the instance $p = \pone$ of $p \psub p = \pzero$
\end{enumerate}
\end{proposition}

\begin{proof}
Each follows by a single rewrite. For example:
\[
p \pmul \pone = \pone \pmul p = p
\]
where the first equality is $\pmul$-commutativity and the second is the $\mathsf{one\_mul}$ axiom.
\end{proof}

\textbf{Axiom count: $48 \to 43$.}

\subsection{De Morgan: Axiom vs.\ Theorem}

The De Morgan laws are currently axioms:
\begin{align}
\pone \psub (a \pmul b) &= (({\pone \psub a}) \padd ({\pone \psub b})) \psub (({\pone \psub a}) \pmul ({\pone \psub b})) \label{eq:dm-and} \\
\pone \psub ((a \padd b) \psub (a \pmul b)) &= ({\pone \psub a}) \pmul ({\pone \psub b}) \label{eq:dm-or}
\end{align}

\begin{proposition}[De Morgan in $[0,1]$]
If $\Prob$ is constructed as the real interval $[0,1]$, then \eqref{eq:dm-and} and \eqref{eq:dm-or} are theorems.
\end{proposition}

\begin{proof}
Standard algebra. For \eqref{eq:dm-and}: expand $(1-a) + (1-b) - (1-a)(1-b)$ to get $1 - ab$.
\end{proof}

\begin{remark}
This does \emph{not} derive De Morgan from our current axioms. It shows that \emph{if we construct} $\Prob$ as $[0,1]$, the axioms become theorems. Currently, De Morgan remains axiomatic.
\end{remark}

\textbf{Status: Axioms, pending $\Prob$ construction.}

\subsection{Conditional Expectation: Monotonicity from Positivity}

\begin{proposition}
Monotonicity of conditional expectation:
\[
(\forall x,\, f(x) \ple f'(x)) \Rightarrow \E[f | g] \ple \E[f' | g]
\]
is derivable from additivity plus a \emph{positivity} axiom:
\[
(\forall x,\, \pzero \ple f(x)) \Rightarrow \pzero \ple \E[f | g]
\]
\end{proposition}

\begin{proof}[Sketch]
Let $h = f' - f$ (pointwise). Then $h \geq 0$ and $f' = f + h$. By additivity:
\[
\E[f' | g] = \E[f | g] \padd \E[h | g]
\]
By positivity, $\E[h | g] \geq 0$, so $\E[f' | g] \geq \E[f | g]$.
\end{proof}

\begin{remark}
This trades the monotonicity axiom for a positivity axiom---not a net reduction in count. However, positivity is arguably more fundamental: it asserts that expectation preserves non-negativity, from which monotonicity follows. The value is conceptual clarity, not axiom count.
\end{remark}

\textbf{Status: Reorganization possible, not a net reduction.}

\subsection{Summary of Part 1}

\begin{center}
\begin{tabular}{lcl}
\toprule
Item & Effect & Status \\
\midrule
Commutativity reductions & $-5$ axioms & \textbf{Proven} (Lean) \\
De Morgan laws & $-2$ axioms & Requires $\Prob$ construction \\
CondExp mono $\leftrightarrow$ positivity & $\pm 0$ axioms & Reorganization only \\
\midrule
\textbf{Confirmed reduction} & $48 \to 43$ & \\
\textbf{With construction} & $43 \to 41$ & \\
\bottomrule
\end{tabular}
\end{center}

The confirmed reduction is modest: 5 axioms eliminated by trivial observations. The De Morgan laws would become theorems if $\Prob$ were constructed as $[0,1]$, but this requires the construction itself---the subject of future work.

% ============================================================================
\section{Part 2: Finite Probabilistic Zorn}
\label{sec:finite-zorn}
% ============================================================================

% Content to be added

% ============================================================================
\section{Part 3: Projective Limits and Infinite Zorn}
\label{sec:projective-limits}
% ============================================================================

% Content to be added

% ============================================================================
\section{Part 4: Synthetic Probability Type Theory}
\label{sec:synthetic}
% ============================================================================

% Content to be added

% ============================================================================
\section{Conclusion}
% ============================================================================

% Content to be added

\appendix

\section{Complete Axiom List}
\label{app:axioms}

% Full list of current axioms

\section{Lean Implementation}
\label{app:lean}

% Code excerpts from Synthetic.lean

\bibliographystyle{plain}

\end{document}
